\scanpage{20}
\textcircled{a}
\diagram{\begin{tikzcd}[row sep=large,column sep=large]
& L^{\bullet\,\mathrm{triv}}=\mathbb R\deleted{\unsure{(\Gamma,\pi)}}\Gamma_{\unsure{*}}(W_0)^{\mathrm{triv}} \arrow[dl,dashed] \arrow[r,bend left] & \text{on met dessus op. triviale de $\pi$}\\
{\begin{tikzpicture}[baseline=(current bounding box.center)]\node[draw,ellipse,inner sep=6pt] {$\begin{gathered}[R^\bullet]\\\Vert\\\mathbb R(\Gamma,\pi)_!(\mathring M)\end{gathered}$};\end{tikzpicture}} \arrow[r,dashed] & K^\bullet=\mathbb R(\Gamma,\pi)_{\unsure{*}}(M) \arrow[u,"\rho\ \text{compatible aux cup}"] &
\end{tikzcd}}
\margin{dans catégorie dérivée de $\Lambda[\pi]$}

\textcircled{b} Appliquons $\mathbb R\Gamma^\pi$ à \textcircled{a}, on trouve
\diagram{\begin{tikzcd}[row sep=large,column sep=large]
& {\begin{gathered}L^{\bullet\pi}=\mathbb R\Gamma(W_0\times S^1)\simeq\mathbb R\Gamma(W_0)\overset{\mathrm L}{\otimes}_{\Lambda}\mathbb R\Gamma(S^1)\\\simeq\mathbb R\Gamma(W_0)\times\mathbb R\Gamma(W_0)(-1]\end{gathered}} \arrow[dl,dashed]\\
{\begin{tikzpicture}[baseline=(current bounding box.center)]\node[draw,ellipse,inner sep=6pt] {$\begin{gathered}[R^{\bullet\pi}]\\\Vert\\\mathbb R\Gamma_!(\mathring P)\end{gathered}$};\end{tikzpicture}} \arrow[r,dash,dashed] \arrow[d,equal] & K^{\bullet\pi}=\mathbb R\Gamma(P)=\mathbb R\Gamma(W-W_0) \arrow[u,"(\alpha,\beta)",very thick]\\
\mathbb R\Gamma_!(W-W_0)&
\end{tikzcd}}
\margin{Suite exacte de cohomologie relative pour $(P,\partial P=S^1\times W_0)$}
\diagram{\begin{tikzcd}[column sep=large]
\mathbb R\Gamma(P)\simeq\mathbb R\Gamma(W-W_0) \arrow[r,"\text{compatible aux cup}","\alpha"'] &\mathbb R\Gamma(W_0)\\
\mathbb R\Gamma(P)=\mathbb R\Gamma(W-W_0) \arrow[r,"\beta",very thick] &\mathbb R\Gamma(W_0)(-1)
\end{tikzcd}}
\margin{s'identifie sur la cohomologie au $\partial$ de la suite exacte de \unsure{dévolution} qui relie $H^i(W-W_0)$ et $H^{i-1}(W_0)$ ...}

\textcircled{c} Appliquons foncteur oubli $\omega$ à \textcircled{a}, on trouve
\diagram{\begin{tikzcd}[row sep=large,column sep=large]
&L^\bullet=\mathbb R\Gamma(W_0) \arrow[dl,dashed]\\
{\begin{tikzpicture}[baseline=(current bounding box.center)]\node[draw,ellipse,inner sep=6pt] {$\begin{gathered}[R^\omega]\\\Vert\\\mathbb R\Gamma_!(\mathring M)\end{gathered}$};\end{tikzpicture}} \arrow[r,dash,dashed] & K^{\bullet\omega}=\mathbb R\Gamma(M) \arrow[u,"\text{compatible aux cup}"]
\end{tikzcd}}

\textcircled{d}
\margin{NB. c'est le $\beta$ de \textcircled{b}}
\diagram{\begin{tikzcd}[row sep=large,column sep=large]
&\mathbb R\Gamma(P)\simeq\mathbb R\Gamma(W-W_0) \arrow[dl,"\beta"',very thick]\\
\mathbb R\Gamma(W_0)(-2] \arrow[r] & [\mathbb R\Gamma_!(W)] \arrow[u,"\rho\ \text{compatible avec cup}"']
\end{tikzcd}}
[Suite exacte en \unsure{Cohomologie|cohomologie} relative pour $(W,W-W_0)$, ou au choix pour $(W,P)$]
\margin{Permet de trouver $\mathbb R\Gamma(W)$ en fonction de \textcircled{a}}
\mspage{19}
